\usepackage[margin=1in]{geometry}
\usepackage{iftex}
\ifPDFTeX
  \usepackage[T1]{fontenc}
  \usepackage{newtxtext}
  \usepackage{newtxmath}
\else
  \usepackage{fontspec}
  \IfFontExistsTF{Times New Roman}
    {\setmainfont{Times New Roman}[Ligatures=TeX]}
    {\setmainfont{TeX Gyre Termes}[Ligatures=TeX]}
\fi
\usepackage{amsmath}
\usepackage{amssymb}
\usepackage{graphicx}
\usepackage{booktabs}
\usepackage{threeparttable}
\usepackage{tabularx}
\usepackage{adjustbox}
\usepackage{ragged2e}
\usepackage{caption}
\usepackage{subcaption}
\usepackage{natbib}
\usepackage[hidelinks]{hyperref}

% Default manuscript layout: 12 pt class option, Times-style font, left-aligned text.
\AtBeginDocument{\RaggedRight}

% Left-align the title block instead of the article class default centered title.
\makeatletter
\renewcommand{\maketitle}{%
  \begin{flushleft}
    {\LARGE\bfseries \@title\par}
    \vskip 1em
    {\large \@author\par}
    \vskip 1em
    {\@date\par}
  \end{flushleft}
}
\makeatother

% Optional local three-line table style; follow journal instructions instead when they differ.
% Top and bottom rules are 1.5 pt; header/body and grouped-header rules are 0.5 pt.
\setlength{\heavyrulewidth}{1.5pt}
\setlength{\lightrulewidth}{0.5pt}
\setlength{\cmidrulewidth}{0.5pt}

\captionsetup{font=small, labelfont=bf, justification=raggedright, singlelinecheck=false}

% Keep table notes compact and readable.
\newcommand{\tablenote}[1]{\begin{tablenotes}[flushleft]\footnotesize\item #1\end{tablenotes}}
